\chapter{Discussion}
\label{chap:discussion}

This chapter contextualises the evaluation findings within the broader literature, reflects on what the novel contributions do and do not achieve, and considers where the work could go next. I try to be honest about the results, including where they were weaker than I expected.

\section{Against the Research Objectives}
\label{sec:against_objectives}

Before discussing specific findings in detail, it is worth stepping back and asking honestly whether the project delivered what Chapter~\ref{chap:introduction} promised. The stated aim (Section~\ref{sec:novel_contributions}) was to deliver three primary novel contributions: a \textit{finance-emotion taxonomy}, a \textit{lightweight ABSA module}, and \textit{end-to-end XAI integration}. These were underpinned by two supporting contributions, namely multi-source data aggregation and a statistical sentiment-price correlation engine. I assess each in turn below, keeping the novelty framing identical to that used in the Abstract and Conclusion.

\textbf{Novelty 1, Finance-emotion taxonomy.} The six-class taxonomy (fear, optimism, uncertainty, confidence, scepticism, mixed) is implemented and integrated into the enrichment pipeline, combining FinBERT probabilities, normalised Shannon entropy, domain lexicons, and aspect priors within a single inference pass. The module is deployed in production and annotates all stored records. A component-contribution ablation (Chapter~\ref{chap:evaluation}) quantifies which signals genuinely affect dominant-label selection, but a full quantitative accuracy evaluation of the taxonomy itself requires an independently labelled emotion-level dataset that does not yet exist for this domain. I consider the taxonomy deployed and partially validated rather than fully validated.

\textbf{Novelty 2, Lightweight ABSA.} The ABSA-lite module extracts aspect-tagged evidence across six financial aspect categories (revenue, growth, risk, valuation, product/AI, macro/policy) using a keyword-triggered, sentence-level design that avoids the computational cost and labelled-data requirement of full sequence-model ABSA \citep{pontiki2014semeval}. It is integrated with the sentiment pipeline and attaches evidence snippets to each stored record. This was delivered in full and is in continuous use on the live corpus.

\textbf{Novelty 3, End-to-end XAI integration.} LIME explanations are served from a live FastAPI endpoint and rendered as token-level heatmaps in the React frontend. Deploying XAI through the full application stack rather than confining it to an offline notebook is the aspect of the system most directly novel relative to \citeauthor{yeo2023xai}'s (\citeyear{yeo2023xai}) review, and it was delivered in full, including the async/CPU-blocking resolution via \texttt{run\_in\_executor}.

\textbf{Supporting contribution, Multi-source aggregation.} The system collects and normalises data from Reddit, X/Twitter, and financial news via three independent ingestion pipelines, storing labelled records in a shared PostgreSQL schema. This was achieved as specified, with an honest caveat: the X API v2 rate and quota limits bounded the Twitter subcorpus below the volume obtainable from Reddit and NewsAPI, which I discuss in Section~\ref{sec:twitter_reflection} and Section~\ref{sec:correlation_discussion}.

\textbf{Supporting contribution, Correlation engine.} The engine implements Pearson, Spearman, rolling correlation, configurable lag analysis, and Granger causality within a single interactive interface. All five methods are functional and user-configurable. The results were more heterogeneous than I anticipated, which I discuss in Section~\ref{sec:correlation_discussion}.

Overall, all three novel contributions were delivered, with the emotion taxonomy partially validated rather than fully validated, and both supporting contributions were delivered as specified. This feels like an honest assessment: the contributions work, but claiming full validation of the taxonomy without a purpose-built emotion-annotated corpus would overstate the case.

\section{Interpreting the Classification Results}

FinBERT's 82.0\% accuracy and 80.3\% macro F1 on my 200-sample benchmark are consistent with the performance ceiling reported for zero-shot FinBERT in the literature. \citeauthor{araci2019finbert} (\citeyear{araci2019finbert}) reported around 86\% on the Financial PhraseBank~\citep{malo2014financialphrasebank} with three-way annotation agreement, a setting that deliberately excludes ambiguous cases. My slightly lower result is what I would expect, given that my dataset includes the noisy, informal social media text the deployed system actually encounters. The 13.5 percentage-point advantage over the keyword baseline mirrors improvements found in prior comparative studies \citep{malo2014financialphrasebank}, which I found reassuring as external validation.

The negative-class recall jump (47.1\% $\to$ 94.3\%) is the one I would show first in a presentation because it is easy to explain to a non-specialist.  The keyword baseline's poor recall on negative samples reflects a fundamental limitation of lexicon-based methods: they depend on explicit negative vocabulary (``loss'', ``decline'', ``downgrade'') but fail on implicit negativity expressed through context-dependent phrasing, for example ``margins came under pressure'', ``guidance proved optimistic'', or ``fell short of elevated expectations''.  FinBERT's bidirectional contextual encoding handles these cases correctly because its pre-training exposed it to the distributional patterns of such phrases in financial corpora.

The neutral class asymmetry (FinBERT: P=97.0\%, R=53.3\%) warrants careful interpretation.  A high-precision, low-recall neutral classifier means the system is conservative: it labels as neutral only when it is very confident that no sentiment-bearing signal is present, defaulting to positive or negative otherwise.  For a financial dashboard, this conservatism may be appropriate. Erroneously classifying a strongly positive announcement as neutral would be more misleading to users than classifying a borderline neutral statement as mildly positive.  Nevertheless, the 28/60 neutral misclassification rate represents a non-trivial source of label noise in the database, particularly for content from financial news that is predominantly factual.  The normalised entropy score produced by the enrichment module provides a partial remedy. High-entropy predictions flag texts where FinBERT is uncertain, allowing the frontend to indicate to users that a label should be interpreted cautiously.

\section{Novel: Finance-Emotion Taxonomy -- Novelty and Implications}

The finance-emotion taxonomy is the most conceptually novel component of this system.  By integrating FinBERT's probabilistic outputs with Shannon entropy, domain lexicons, and aspect priors within a unified inference function, the taxonomy translates a coarse three-class sentiment label into a behaviourally richer six-class representation that more closely mirrors the vocabulary used by financial commentators and practitioners.

The distinction between \textit{fear} and \textit{scepticism}, for instance, is meaningful in a behavioural finance context: fear is associated with immediate, reactive selling behaviour, while scepticism reflects a sustained bearish outlook that resists positive news \citep{bollen2011twitter}.  Similarly, \textit{uncertainty}, detected primarily through elevated softmax entropy rather than lexical signals, captures the market-moving potential of ambiguity itself, consistent with Tetlock's \citeyear{tetlock2007giving} finding that negative media language can drive trading even in the absence of new information.

The taxonomy is deployed as part of the delivered system, but its accuracy against an independently emotion-labelled corpus has not been measured in this dissertation because no such corpus exists in the public domain for the finance emotion classes used here. Constructing that corpus and running a formal accuracy evaluation are identified as future work in Section~\ref{sec:future_work}, where a suitable publication pathway is also briefly discussed.

\section{Novel: ABSA-Lite -- Lightweight Aspect Extraction}

The ABSA-lite module addresses the gap identified in the literature between full sequence-model ABSA \citep{pontiki2014semeval} and the absence of any aspect analysis.  Its keyword-triggered, sentence-level evidence extraction avoids the computational cost and labelled-data requirement of fine-tuned ABSA systems, making it suitable for continuous, high-throughput ingestion pipelines.

The practical value is twofold.  For end users of the dashboard, the evidence snippets provide a transparent audit trail: rather than seeing only a ``negative'' label for a company, users can see the specific sentence that triggered the \textit{risk\_liquidity} aspect, providing a level of granularity that directly supports the explainability objectives of the system.  For researchers, the aspect distribution across the 150{,}000-record corpus provides a lens on \textit{what investors are discussing}, not just \textit{how they feel} about it, a dimension underexplored in the prior literature.

\section{Novel: XAI in Production -- Closing the Deployment Gap}

The integration of LIME into the FastAPI REST endpoint and the React SPA frontend directly addresses the gap identified by \citeauthor{yeo2023xai} (\citeyear{yeo2023xai}): that most XAI applications in finance remain confined to offline tools rather than production user interfaces.  Three architectural decisions together demonstrate the engineering considerations required to bridge research-grade XAI tooling and production deployment: dispatching the LIME perturbation sampler onto the default thread-pool executor via \texttt{asyncio.to\_thread} (the Python~3.9+ wrapper over \texttt{loop.run\_in\_executor}), allocating a dedicated 120-second frontend timeout, and returning token-level heatmaps as structured JSON.

The LIME evaluation in Chapter~\ref{chap:evaluation} produced one particularly instructive finding that deserves emphasis here.  Running LIME on the sentence ``The Federal Reserve kept interest rates unchanged, maintaining its cautious outlook on inflation'' revealed that FinBERT predicted \textit{negative} (55.9\% confidence), with LIME identifying ``cautious'' and ``inflation'' as the principal drivers.  Both are legitimate negative signals in many financial contexts, but in this sentence they are neutral-framed; the model has no mechanism to distinguish the communicative intent of caution-as-prudence from caution-as-concern.

This is not just a curiosity: it directly illustrates \textit{why} the neutral recall deficit observed in the confusion matrix exists.  Without LIME, the neutral recall number (53.3\%) is a symptom without a diagnosis.  With it, the specific failure mode is identifiable: contextually neutral risk vocabulary systematically pushes FinBERT toward the negative class.  One targeted intervention, for instance fine-tuning on a corpus of central bank communications where cautious language is categorically neutral, would likely improve neutral recall substantially without degrading performance elsewhere.  This is exactly the kind of insight that justifies embedding XAI within a diagnostic toolchain rather than treating it purely as a user-facing feature.

The decision to offer SHAP for offline auditing rather than real-time requests reflects a principled trade-off consistent with the literature: \citeauthor{lundberg2017shap} (\citeyear{lundberg2017shap})'s Shapley values provide stronger theoretical guarantees than LIME's linear surrogate, but their computational cost renders them unsuitable for interactive query patterns.  Future work could explore kernelSHAP with reduced evaluation budgets or distilled explanation models to bring theoretically grounded explanations closer to real-time feasibility.

\section{Sentiment-Price Correlation: Observations and Caution}
\label{sec:correlation_discussion}

Combining Pearson and Spearman correlation, configurable lag analysis, Granger causality, rolling correlation, and market adjustment within a single interactive interface goes beyond what any single prior study in the reviewed literature had implemented. The ability to select a ticker, adjust the window, and compare multiple statistical methods simultaneously supports a level of exploratory analysis that existing tools do not currently offer.

That said, the results were more modest than initial expectations suggested. Studies such as Bollen et al.\ \citeyear{bollen2011twitter} reported relatively consistent correlations between Twitter mood and market movements, which led to an expectation of similar consistency across tickers. In practice, Granger causality tests were statistically significant for some tickers but not others, correlation strength varied substantially with the choice of window, and macro market events appeared to drive both sentiment and price simultaneously, making causal attribution difficult. This is not a failure of the system; it reflects the genuine complexity of the underlying relationships. The dashboard surfaces that complexity transparently rather than smoothing it into a misleadingly clean output.

The cautious interpretation of all these results is essential, and I embedded that caution prominently in the Methodology and Legal pages of the dashboard. Several well-documented challenges limit what conclusions can be drawn:

\begin{itemize}
    \item \textbf{Non-stationarity.}  Both sentiment and price return series are typically non-stationary over long windows, which can inflate correlation statistics.  The implementation performs augmented Dickey-Fuller tests before Granger analysis as a precaution, but daily windows may still exhibit non-trivial drift.
    \item \textbf{Confounding.}  Macro events (earnings seasons, rate decisions) simultaneously drive both sentiment and price, creating confounded co-movement that does not reflect a causal sentiment $\to$ price channel.
    \item \textbf{Sample sizes.}  For individual tickers with fewer than 20 to 30 daily data points in the selected window, correlation estimates are highly unstable; the data quality endpoint flags this risk for users.
    \item \textbf{Macro regime.}  A recurring theme across the expert interviews (Appendix~\ref{app:interviews}) was that post-2008 quantitative easing and the ensuing asset-price inflation have partially decoupled market prices from company-level fundamentals, weakening purely technical signals relative to earlier market regimes. This matters for interpreting the correlation results. In a regime where price movements are partly driven by macro liquidity conditions rather than by information flow, even a faithfully implemented sentiment-price correlation engine is measuring a relationship whose signal-to-noise ratio varies with the wider economic backdrop. This is not an argument against running the analysis. It is an argument for presenting the outputs with the window-dependent and regime-dependent caveats that the dashboard already surfaces.
\end{itemize}

These caveats notwithstanding, the aggregated multi-source corpus (over 150{,}000 records) represents a substantially larger evidence base than most single-study datasets, providing reasonable statistical power for tickers with sufficient daily mention volume over extended periods.

\section{Multi-Source Aggregation and Ingestion Resilience}

A recurring theme in the sentiment analysis literature is the reliance on a single data source.
Tetlock's foundational study \citep{tetlock2007giving} used only \textit{Wall Street Journal} columns; Bollen et al.\ \citep{bollen2011twitter} used only Twitter; Kraaijeveld and De Smedt \citep{kraaijeveld2020predictive} used only Twitter for cryptocurrency markets.
Each approach captures a different slice of the information landscape, but none accounts for the divergence that regularly occurs between these channels.
Institutional news tends to react to verified, price-relevant events; Twitter and Reddit often lead with speculation, retail momentum, and meme-driven behaviour that is structurally different in both timing and reliability.

By collecting from all three sources and tagging each record with its origin, the system allows queries that explicitly compare sentiment across channels rather than merging them into a single aggregate.
Long et al.\ \citep{long2023reddit} showed that Reddit's r/WallStreetBets sentiment during the GameStop rally was largely disconnected from broader market signals, suggesting that platform-separated analysis captures dynamics that cross-source aggregation would suppress.
The correlation engine's per-source filtering is therefore not a convenience feature; it reflects a genuine analytical choice.

\subsection{X/Twitter: Working Within Platform Limits}
\label{sec:twitter_reflection}

The Twitter pipeline is deliberately restricted to the official X API v2 to remain within the platform's terms of service and the project's ethics commitments (Appendix~\ref{app:ethics}). The practical consequence of that decision is the single biggest honest limitation of this project's data layer. Free-tier and low-tier access to the v2 endpoints caps both monthly tweet volume and the depth of historical queries. These limits were substantially tightened during the 2023--2024 platform changes, which also raised the price of elevated access well beyond an undergraduate project budget. Within those constraints, and even with efficient query construction and exponential backoff on 429 responses, the Twitter subcorpus is smaller and less temporally deep than the Reddit and NewsAPI subcorpora. That imbalance is visible downstream. Per-ticker correlation windows that include Twitter sentiment carry fewer daily data points, and Twitter-only views have correspondingly lower statistical power than Reddit- or news-only views. I therefore treat the Twitter signal as a secondary, corroborating input rather than a primary one. The per-source filters in the dashboard make this imbalance transparent rather than hiding it inside an aggregate. A historical CSV path is retained alongside the live API to enable fully reproducible evaluation against a fixed corpus, which single-pathway designs cannot offer.

\section{Reflections on the Build Process}

Two points in the build process proved particularly instructive. The first was an extended debugging session on the NewsAPI deduplication logic: database row counts appeared substantially higher than expected, and systematic inspection of content hashes revealed that syndicated wire copy was repeating under different publisher names. Identifying and handling that duplication was time-consuming but necessary for the integrity of the downstream charts. The second was the first successful LIME run on a sentence selected because it was intuitively positive: the token weights aligned with the words that a domain reader would identify as sentiment-bearing (``surged'', ``guidance''). That result provided early confidence in the explainability integration as a meaningful analytical tool rather than a post-hoc visual artefact.

\section{Prototype Polish and Commercial Context}

The screenshots in Section~\ref{sec:frontend_screenshots} illustrate that the system extends beyond a backend data pipeline: the React SPA provides multi-source views, per-ticker drill-down, explainability, and correlation within a single interface. This is relevant for demonstrating the practical utility of the work to potential investors or product teams, though it is distinct from claiming the build constitutes a commercial offering. A real rollout would need hardened authentication, billing, observability, contractual data licences, model governance sign-off, and jurisdiction-specific compliance (for example clear boundaries around investment advice in the UK). I kept the Legal and Methodology routes in the SPA precisely so those gaps are visible rather than hidden.

\section{Ethical Reflection}

The system's adherence to platform terms of service and its anonymisation of user data reflect an appropriate approach to the ethical use of publicly available social media content.  However, the deployment of a sentiment analysis system that aggregates and analyses investor opinions raises broader questions about information asymmetry: dashboard users with access to real-time sentiment aggregates and correlation analytics have an informational advantage over those without.  This connects to a long-standing debate in financial economics about whether markets are informationally efficient \citep{fama1970efficient} and whether sentiment-driven traders introduce persistent mispricings \citep{delong1990noise, barberis2003behavioral}.  While the present system is designed for academic purposes, commercial deployment would require careful consideration of fairness, accessibility, and potential regulatory implications under financial services legislation.

The Groq-backed AI narrative feature, which generates natural-language summaries of ticker sentiment grounded in database records, introduces an additional layer of consideration. LLM-generated text may present statistical artefacts as confident narratives. That risk is mitigated by grounding the narrative in database counts and explicitly disclaiming non-advisory intent on the Legal page.  The broader principle is that model outputs should be accompanied by clear documentation of their intended use, limitations, and potential failure modes. This principle is formalised in the Model Cards framework of Mitchell et al.\ \citep{mitchell2019modelcards}, which the Methodology page of the dashboard is designed to approximate.
